\chapter{Persistent agent state, durable workflows, and idempotent retries}\label{persistent-agent-state-durable-workflows-and-idempotent-retries}

Netflix once lost roughly 4 percent of cloud-operations deployments to transient failures, even though the service already contained homegrown retries and compensating-transaction code. Meyers and Zienert (\href{https://netflixtechblog.com/how-temporal-powers-reliable-cloud-operations-at-netflix-73c69ccb5953}{2025}) reported that the team moved coordination into a stateful service that records progress and reissues work after failures. This durable workflow engine cut the reported rate to 0.0001 percent, and the team deleted the orchestration code it had replaced. The account comes from the adopting team. It is vendor-adjacent and was not independently audited, but it is unusually concrete about both the migration and the observed change.

That account is this chapter's only strong item. Historical surveys, systems papers, self-evaluated systems, and practitioner reports supply the remaining evidence, and controlled experiments on agents are absent. I therefore make a narrow design argument and prescribe explicit artifacts, boundaries, and failure checks. None of the prescriptions depends on a numerical threshold.

An agent that dies after step nine of a fourteen-step run presents the same coordination problem at smaller scale. The process has disappeared, while its completed work and its external effects may remain. \textbf{Durable execution} is a run's ability to survive process death and resume from recorded progress without duplicating completed effects. That property requires explicit state, a coordinator that owns recovery once coordination becomes complex, and retryable steps whose repeated execution does not repeat their effects.

The hard case is a process that changes the world and dies before recording what it changed. A plan left only in model context, a tool result held only in memory, or a completion flag scheduled as the next write all disappear at the point when recovery needs them. The surviving system then has to separate completed work, incomplete work, and work whose status cannot be determined. Without that separation, the restart presents the old run identity without its evidence.

\section{Declare the state that outlives the worker}\label{declare-the-state-that-outlives-the-worker}

Consider the step-nine failure with no completion record on disk. The first eight results may exist only in the lost context, while the ninth may already have changed an external system. The support for this section is historical inference from stream-processing research, self-evaluated systems, and practitioner accounts. No controlled experiment establishes the result for agents.

Persisted run state is the substrate recovery reads. The evidence does not establish how well this design works across open-ended agent workloads. I therefore treat explicit state as an architectural boundary to test rather than a measured guarantee that recovery succeeds.

The first decision is what counts as control state. The current plan records intended work, and progress records the subset accepted as complete. A plan changes through explicit revisions. Progress advances through named steps and carries the evidence required to accept each transition.

The run also accumulates evidentiary state: intermediate knowledge, results returned by tools, and session events. Knowledge snapshots replace older versions under a declared rule. Tool results belong to specific invocations, and events append in order. One opaque transcript preserves their text while obscuring their different identities and update rules.

Each recoverable item should be a declared artifact with a stable identity, an owner, and an update rule. A plan can carry a run identity and a revision number. A tool outcome can belong to a step identity plus an invocation identity. A progress record can name the last completed step and the evidence that completion produced. An event record can append the attempted action, its observed result, and the state transition that followed. Stable identities let a replacement worker refer to the same work instead of creating a parallel copy.

Osmani (\href{https://addyo.substack.com/p/long-running-agents}{2026}) recommended three such artifacts kept outside the model context: a plan file, progress notes, and an append-only event log. The same arrangement allows a full context reset to be rebuilt from a handoff file. That is a practitioner recommendation rather than a measured result. A progress record the agent writes about its own work has a further problem. It is a self-report generated from the same context that may already be wrong. Completion evidence should therefore come from the effect boundary, e.g., a tool response, a commit identifier, or an independent state read, rather than from the agent's assertion that a step finished.

Two stores that both appear authoritative create a new failure mode. If the queue says step nine is ready while a progress file says it completed, recovery depends on an undocumented precedence rule. A declared design names the source of truth for each fact and treats other copies as indexes or caches. The queue may own delivery state, the event log may own execution history, and a versioned snapshot may accelerate reconstruction without overriding later events.

This architecture assumes the worker is disposable. A fresh worker reads the durable queue record, loads the latest valid snapshot, applies later events, and reconstructs the next allowed action. It does not need the failed process's heap or model context. The agent is amnesiac. The storage layer is responsible for remembering enough that the amnesia does not break the evidence chain.

That division resembles a change that took roughly two decades in stream-processing systems. Fragkoulis et al.~(\href{https://arxiv.org/abs/2008.00842}{2020}) described early systems that treated state as application-managed data and later systems that brought state, checkpoints, and recovery under runtime control. The survey supplies directional evidence and ran no experiment on agents. Its useful inference is architectural. Recovery became state-centric once runtimes could identify the state they had to preserve and coordinate it with progress through the input.

The analogy has a limit. A stream processor usually applies a specified computation to structured input. An agent may revise its plan, interpret ambiguous evidence, or call a model that returns a different answer to the same prompt. Explicit state cannot make those decisions deterministic. It can preserve which decision occurred, which evidence informed it, and which work followed, so a replacement worker does not silently substitute a new history.

The ADEMA architecture, from Hanlin (Zhou) et al.~(\href{https://arxiv.org/abs/2604.25849}{2026}), is a self-evaluated system that offers a narrower check on this reasoning. In a fixed matrix of 60 runs, removing checkpoint and resume produced the only invalid run. That result isolates checkpointing inside one system and one bounded experiment. It does not compare recovery designs or establish a general failure rate. The study showed that a knowledge-state snapshot can preserve information across an interruption when the system has defined what belongs in the snapshot.

Halukurike et al.~(\href{https://tech.instacart.com/blueberry-force-multiplier-for-the-on-call-engineer-98c446dfcc12}{2026}) reported a database-backed durable queue with per-pass session snapshots handling approximately 25,000 diagnostic passes per month at 99.9 percent success, including takeover by another worker during a pass. The team that built the system reported those numbers, and no independent evaluation is on record. The workload also consists of bounded diagnostic passes. A multi-day agent that keeps discovering new subproblems may accumulate state whose relevance, size, and ownership change during the run.

Longer work therefore needs more than periodic serialization. Each snapshot should state which event position it incorporates, which software or schema version can read it, and which prior artifacts remain authoritative. A replacement worker must reject a snapshot it cannot interpret, because partial loading invents state.

Retention is a separate design decision. Keeping every model response, tool payload, and intermediate artifact indefinitely turns recovery storage into an uncontrolled cost and an uncontrolled audit surface. A retention rule must preserve the evidence needed for recovery while defining when older material can be compacted or removed. The deduplication records described later in this chapter must be preserved under that rule, because deleting the record of a completed invocation restores the duplicate-execution risk the record existed to remove.

The restore path needs its own test. Start a run, let it produce a plan, a tool result, a progress update, and an event sequence, then remove the worker. A new worker should reconstruct the same completed-step set and choose the same next eligible step from the declared artifacts alone. Any reliance on a live session, an untracked local file, or an operator's memory identifies state that the architecture has not yet declared.

I learned that distinction from a contrary case in my own estate. A 2026 audit of my contribution pipeline found that all thirteen workflow skills constituting the pipeline existed only as unversioned local files, with zero backup. That is an uncomfortable finding for the author of a chapter about declared state. The files were persistent in the weak sense that they survived process exit, but they had no version history and no tested restore path. I moved them into a versioned repository, because a declared artifact that cannot be restored is still an operational single point of failure.

Persistence also leaves several agent problems untouched. It does not bound model cost, prevent a long run from drifting away from its original objective, or remove the work of auditing accumulated decisions. In some systems it can make those burdens larger by preserving every questionable intermediate judgment. The narrower benefit is recoverability. After process death, the next worker can determine what the system knew, what it recorded as complete, and where uncertainty begins.

\section{When coordination belongs in an engine}\label{when-coordination-belongs-in-an-engine}

The deleted orchestration code at Netflix is as informative as the reported failure reduction. Before the migration, application code coordinated retries, compensating actions, and service state. Afterward, the workflow engine held the execution history and scheduled application operations from that history. The strong practitioner account supports the result for this cloud-operations service. It does not isolate which part of the migration caused the change, and it does not establish the same effect for agent workloads.

An engine is worth its cost when it owns facts that no individual worker can know reliably after a crash. It records which step became eligible, which attempt started, which result committed, which external event arrived, and which retry policy applies next. Workers still perform model calls, tool calls, and domain operations. The engine owns their order and lifecycle. A replacement worker therefore receives work from recorded history instead of reconstructing history from local conventions.

This division places cross-service consistency in one architectural layer. Consider a deployment that reserves capacity, changes traffic, and updates an inventory. Each service can make its local transaction correct while the deployment as a whole remains half-finished. A saga coordinates those local transactions as one logical unit and specifies a compensating action for each step that cannot be rolled back directly. When every service encodes its own retry and compensation conventions, no component holds a complete view of the deployment.

The workflow engine makes that global view durable. A failed capacity reservation remains visible with its retry policy. A completed traffic change stays recorded while the engine waits for a compensating operation. The operator can inspect one execution history instead of inferring order from several service logs. Failure localization improves because the coordinator records the control path and the workers report their results into it.

The supporting literature is useful but limited. Laigner et al.~(\href{https://arxiv.org/abs/2103.00170}{2021}) combined a literature review, repository analysis, and a survey of more than 120 practitioners, and found reliability problems concentrated around hand-built sagas and convention-managed consistency. Nadeem and Malik (\href{https://arxiv.org/abs/2204.07210}{2022}) ported a benchmark system with 22 known bugs to a workflow engine in a single-participant case study and reported that the bugs became easier to localize. Their debugging times were compared with times reported separately in another study, so they do not form a controlled head-to-head measurement.

Centralized history does not supply a universally correct consistency guarantee. Zhang et al.~(\href{https://arxiv.org/abs/2208.09827}{2022}) surveyed a decade of transactional stream-processing work and found no generally accepted approach, for computations far more deterministic than agent workflows. Each system chose guarantees around particular application characteristics. That null result is load-bearing for the adoption decision. An engine cannot be selected by checking a feature box labeled `exactly once' and assuming that the application's state, latency, and effects now fit the guarantee.

The engine boundary should follow the workload. I use three questions as a design check. Does meaningful work remain exposed if the process dies mid-flight? Does the workflow wait on external events? Does it perform irreversible external side effects? A yes to any of them identifies state that a stateless scheduler cannot reconstruct by reading the current source of record. Three noes usually point toward a timer, a lockfile, and a fresh read from that source.

That three-question check is my operational rule, and no experiment produced it as a threshold. In one \textbf{unverified working artifact} from my own live maintenance loop, each run performed approximately 44 seconds of work during a 120-minute window. An overlap-skip guard fired zero times across 77 runs. An unverified working artifact has a source that is uncommitted at the stated repository revision, with figures read out of that repository rather than independently re-measured. It can show the mechanism of a single system, but it cannot establish a rate or support a general claim.

That loop had no valuable mid-flight state, no wait on an external event, and no irreversible effect. A workflow engine would therefore have added history and deployment machinery without improving the measured behavior.

The opposite workload has a different shape. An agent may wait overnight for approval, call several services that cannot share a transaction, or spend substantial money on a model result that should survive a worker restart. A process-local retry loop loses its authority when the process dies. A workflow engine can persist the wait, record the expensive result, and assign the next operation to another worker while preserving one execution identity.

That capability constrains implementation. Engines that reconstruct control flow from recorded history require workflow code to make the same orchestration decisions when it reads the same history. A clock read, a random branch, or a changed iteration order inside that code can select a path absent from the existing history. Code evolution therefore needs version boundaries. Even a change from a fixed workflow signature to variable arguments, which looks harmless, can make a long-running execution incompatible with its recorded input.

Decomposition creates costs of its own. Child workflows used only to organize code add another execution to inspect and another history boundary to understand. Central coordination can reduce throughput or local autonomy when every small action must pass through one scheduler and one persistence layer. The engine improves visibility by concentrating control, and concentrating control also concentrates contention and operational responsibility.

The integration boundary often costs more than the workflow definition. My own durable media integration, an \textbf{unverified working artifact}, took 15 commits across 45 files and added 11,662 lines, including 4,665 test lines. The workflow definition occupied a small fraction of that change. Most of the work established clean payloads, stable execution identities, idempotent external requests, injected-failure gates, and reconciliation with durable storage. These counts illustrate one system and support no industry cost estimate.

An engine should therefore own coordination without absorbing domain behavior. Application code still decides what a valid deployment, payment, or agent result means. The platform owns durable history, retry scheduling, waiting, and the transition between recorded step states. Keeping that boundary explicit makes engine replacement possible and keeps application correctness from depending on undocumented behavior in a particular worker.

\section{The interval that cannot be closed}\label{the-interval-that-cannot-be-closed}

A worker sends a merge request, receives success, and dies before it can record the completed step. The engine sees an unfinished attempt and sends the step again. My treatment of this interval rests on directional systems research and practitioner accounts. There is no controlled result showing that one idempotency design makes agent workflows reliable across workloads.

The redelivery is correct, because the engine cannot infer completion from a missing record. This policy is at-least-once delivery. The runtime keeps delivering work until it holds a recorded completion, and one logical step may therefore receive several execution attempts. Durable execution can record a committed step result once in its own history. It cannot close the interval between an external system accepting the operation and the worker committing that result to the history.

That interval remains even when both systems are individually reliable. Suppose attempt A begins step nine under invocation identity \texttt{run-42/step-9}. The worker asks a code host to merge a change, and the host commits the merge. The process stops before the worker records the returned result. The engine schedules attempt B, because its last durable fact says that step nine started. The code host has already moved to a new state, while the workflow history still describes an in-flight operation.

The invocation identity has to cross that boundary. Attempt B must send the same idempotency key that attempt A sent. Morling (\href{https://www.morling.dev/blog/building-durable-execution-engine-with-sqlite/}{2025}) described idempotency keys at side-effect boundaries as what makes the execute-then-log crash window safe on replay. The code host then either returns the stored response for that key or recognizes that the requested state already exists and reports the same observable result. An operation is idempotent when applying it repeatedly converges on the same observable state as applying it once. The useful contract is that a duplicate becomes indistinguishable from no duplicate.

\begin{verbatim}
runtime sends request
    -> external system may commit an effect
        -> runtime durably records completion
           ^                          ^
           |<-- execute-then-log gap->|

process dies inside the gap
    -> status uncertain
        -> ordinary retryable failure
            -> duplicated effects

request known not to have executed
    -> retry

recovery design
    -> distinguish known not executed from status uncertain
\end{verbatim}

Key design determines what the contract covers. A fresh attempt identifier fails, because the downstream system sees attempt B as new work. A key scoped only to the user can collapse two legitimate operations into one. The stable identity should name the logical invocation and include enough version or input identity to distinguish work that is genuinely different. Retries reuse it. New logical work receives a new one.

The downstream implementation also needs an atomic claim on that key. Two workers can receive the same step concurrently after a timeout or a lease dispute. If both check for a cached result and then perform the effect before either writes the cache, the check has added latency without preventing duplication. A unique transactional record, or an equivalent compare-and-set operation, must decide which worker owns the invocation before the effect occurs.

Execution caching can extend this rule through a call graph. Psarakis et al.~(\href{https://arxiv.org/abs/2312.06893}{2023}) described a runtime that caches results by invocation identity, so that a repeated parent call reuses completed child calls. Within that runtime's transactional boundary, retried failed calls, recorded completed results, and call ordering can compose into a strong guarantee. An external service that ignores the invocation identity remains outside the boundary, whatever the runtime calls its delivery semantics.

This scope rules out the blanket claim that a completed invocation does not run twice. A step that succeeded externally and remained unrecorded may run again, because the runtime has no completion to consult. Execution occurs at least once, while results committed inside the runtime boundary are recorded exactly once. Safety for the in-flight external effect still depends on the idempotency contract at that effect's boundary.

Some external systems provide that contract directly. Payment APIs often accept a client-supplied key and bind it to the first accepted request. A database can commit the business change and the invocation record in one transaction. A content-addressed object store can treat a write of identical content under the same name as convergence. When a tool offers none of these, the workflow must either add an adapter that owns deduplication or treat recovery as an unknown-state decision.

Meyers and Zienert (2025) also reported that forcing operations through an idempotency rewrite uncovered and fixed pre-existing retry defects. The rewrite changed application behavior at the engine boundary, and the engine could not infer that contract from the old code. This is directional evidence from the adopting team.

Unknown state requires a safe terminal path. In an \textbf{unverified working artifact} that contains my guarded-mutation recipe, the worker records an intent claim before an irreversible action, performs the action, and then resolves the claim with the observed result. A resolved claim can return its cached result on redelivery. An unresolved claim after recovery means the effect may or may not have occurred. The workflow therefore stops and asks a human to reconcile the external state.

Guessing in either direction is unsafe. Assumed success loses work that never happened, and assumed failure repeats an irreversible effect.

An \textbf{unverified working artifact} from my own fault demonstration makes that boundary visible. A naive pipeline was killed after it requested a merge and before it recorded completion. Its retry requested the merge again, and the workflow reported success without alerting anyone. Under the same kill placement, the guarded variant requested one merge, emitted one escalation, and marked the workflow failed. The failed outcome was the more reliable one, because it preserved the uncertainty instead of fabricating a completed history.

Idempotence also belongs at the step boundary before the external boundary. Model calls are canonical durable steps, because they can be slow, costly, and irreproducible. Once a model result commits under a stable invocation identity, a retry can reuse it and continue with the same evidence. Reissuing the call may change both cost and control flow even when no external database has changed.

The boundary should not surround every function call. Each recorded step adds scheduling, serialization, storage, and history-reading work. Pure calculations that are cheap and deterministic can run again. A durable boundary is justified when the work is expensive, slow, externally visible, or impossible to reproduce from recorded inputs.

A boundary can also be too broad. In another \textbf{unverified working artifact} from my own system, a transient read sat inside the same deduplication unit as an irreversible mutation. When the read failed, the system stamped the whole unit terminal, poisoned the claim, and went silent about work it had never performed. A watchdog consulted the same fail-closed store and stopped reporting at the crash. This case is narrative illustration and supplies no evidence for the general design claim. Separating the retryable read from the mutation claim would have preserved an accurate pending state.

Trofimov et al.~(\href{https://arxiv.org/abs/1907.06250}{2019}) proposed describing guarantees through determinism and observable outcomes. Equal inputs should yield equal outputs, and a duplicate should leave no observable difference. That formulation directs attention to what callers can test. It remains one research group's proposal and has not displaced at-least-once and exactly-once vocabulary in common system descriptions.

Repeated execution can converge perfectly on the wrong value. The step still needs postconditions that check its domain result, and the workflow still needs compensation or escalation when the external system cannot make the operation idempotent. Retry machinery addresses neither responsibility.

\section{Companion patterns}\label{companion-patterns}

Eleven adjacent entries in the companion catalog refine this recovery contract without adding another architectural decision. They cover external writes that carry receipts, idempotency keys, and compensators, retry behavior expressed as falsifiable invariants, snapshots taken off the critical path at consistency boundaries, and a checkpoint cadence set from measured costs and failure rates. Others cover execution records committed together with the data change, shared agent state placed behind transactions, event histories started with a weak versioned schema, and operations designed to tolerate reordering. Three entries have thin support and provide no evidence for anything claimed in this chapter. They propose writing event-coordination guarantees explicitly, attaching postconditions to stochastic steps, and converting repeated plans into reusable blueprints.

\section{Run the crash check}\label{run-the-crash-check}

Pick the agent workflow whose mid-run failure has the highest consequence. Give the run a stable identity, and write its plan, its progress, and its append-only events as declared artifacts with named owners. Stop the process after several steps and delete nothing. A replacement worker should identify completed work, pending work, and unknown external state from those artifacts alone.

Run the three-question adoption check before introducing an engine. Look for valuable state exposed mid-flight, waits on external events, and irreversible side effects. If all three are absent, read the source of record again under a timer and a lock. If any are present, document which coordination facts the engine will own and which domain decisions remain in application code.

Then select the retried step with the most dangerous external effect. Give the logical invocation a stable idempotency key, propagate it into the side-effecting system, and cache the completed result under the same identity. Kill the worker after the effect succeeds and before the completion record commits. The recovered run should return the prior result, converge on the same external state, or stop with an explicit unknown-state escalation.

Recorded state makes a run resumable. Chapter 9 makes it replayable and puts the recovery claim under measurement.

\section{Sources and evidence}\label{sources-and-evidence}

The evidence class and strength on each entry below come from its catalog record. Author-system cases in this chapter are narrative illustration and are not part of the evidence base.

\subsection{Make agent state a first-class persistent artifact}\label{make-agent-state-a-first-class-persistent-artifact}

\begin{itemize}
\tightlist
\item
  lit/directional: Fragkoulis, Carbone, Kalavri \& Katsifodimos (2020). A Survey on the Evolution of Stream Processing Systems. The VLDB Journal (2024). arXiv:2008.00842. (Exactly-once recovery arrived only when state became a first-class managed runtime artifact; implicit-state systems had lossy recovery.)
\item
  lit/anecdotal: Hanlin (Zhou) et al.~(2026). ADEMA: A Knowledge-State Orchestration Architecture for Long-Horizon Knowledge Synthesis with LLM Agents. arXiv:2604.25849. (In a fixed 60-run matrix, removing checkpoint/resume produced the only invalid run.)
\item
  practitioner/directional: Addy Osmani (2026). ``Long-running Agents'', \href{https://addyo.substack.com/p/long-running-agents}{Elevate newsletter}, 2026-04-30. (Plan file, progress notes, append-only event log outside the context make an agent recoverable and enable full context resets rebuilt from a handoff file.)
\item
  practitioner/directional: Karthik Halukurike et al.~(2026). ``Blueberry: Force Multiplier For The On-Call Engineer'', \href{https://tech.instacart.com/blueberry-force-multiplier-for-the-on-call-engineer-98c446dfcc12}{Instacart tech blog}, 2026-07-14.
\end{itemize}

\subsection{Use a durable workflow engine}\label{use-a-durable-workflow-engine}

\begin{itemize}
\tightlist
\item
  lit/anecdotal: Nadeem \& Malik (2022). A Case for Microservices Orchestration Using Workflow Engines. ICSE-NIER. arXiv:2204.07210.
\item
  lit/directional: Laigner, Zhou, Vaz Salles et al.~(2021). Data Management in Microservices: State of the Practice, Challenges, and Research Directions. PVLDB 14(13). arXiv:2103.00170. (SLR + repo analysis + 120+ practitioner survey: hand-rolled sagas and convention-managed consistency are where reliability failures concentrate.)
\item
  practitioner/strong: Jacob Meyers and Rob Zienert (2025). ``How Temporal Powers Reliable Cloud Operations at Netflix'', \href{https://netflixtechblog.com/how-temporal-powers-reliable-cloud-operations-at-netflix-73c69ccb5953}{Netflix TechBlog}, 2025-12-15.
\item
  lit/null-result: Zhang, Soto \& Markl (2022). A Survey on Transactional Stream Processing. The VLDB Journal. arXiv:2208.09827. (Carried as the negative result behind the engine choice; it is about deterministic data pipelines and says nothing about agents.)
\end{itemize}

\subsection{Make retried steps idempotent}\label{make-retried-steps-idempotent}

\begin{itemize}
\tightlist
\item
  lit/directional: Psarakis et al.~(2023). Styx: Transactional Stateful Functions on Streaming Dataflows. SIGMOD line. arXiv:2312.06893. (Execution caching keyed by invocation identity gives exactly-once composition across call graphs.)
\item
  practitioner/directional: Meyers and Zienert (2025), \href{https://netflixtechblog.com/how-temporal-powers-reliable-cloud-operations-at-netflix-73c69ccb5953}{Netflix TechBlog}, 2025-12-15. (The forced idempotency rewrite alone fixed pre-existing retry issues.)
\item
  practitioner/directional: Gunnar Morling (2025). ``Building a Durable Execution Engine with SQLite'', \href{https://www.morling.dev/blog/building-durable-execution-engine-with-sqlite/}{morling.dev}, 2025-11-20. (Idempotency keys at side-effect boundaries make the execute-then-log crash window safe on replay.)
\item
  lit/directional: Trofimov, Kuralenok, Marshalkin \& Novikov (2019). Delivery, consistency, and determinism: rethinking guarantees in distributed stream processing. arXiv:1907.06250.
\end{itemize}
